% answers_short.tex —— 2015 年·全国Ⅰ卷·理科 · 简短答案速查
% 由 tools/make_short_answers.py 自 answers.tex 的 \gksolhead/\gkanswer 提取，
% 并与 2015/2015_数学真题_全国1理_答案要点.md 逐题比对后生成，勿手改（改源 MD 后重跑）。

\begin{gknotice}
\textbf{使用说明}：本节只给最终答案便于快速核对，完整推导见解析版。答案为编写者独立推导，\textbf{非官方标准答案}。
\end{gknotice}

\gkpart{选择题}
\begin{center}\begin{tabular}{ccccccccccccc}
\toprule
题号 & 1 & 2 & 3 & 4 & 5 & 6 & 7 & 8 & 9 & 10 & 11 & 12 \\
\midrule
答案 & A & D & C & A & A & B & A & D & C & C & B & D \\
\bottomrule
\end{tabular}\end{center}

\gkpart{填空题}
\begin{itemize}[leftmargin=2.6em]
  \item[\textbf{13.}] $1$
  \item[\textbf{14.}] $(x-\frac32)^2+y^2=\frac{25}{4}$
  \item[\textbf{15.}] $3$
  \item[\textbf{16.}] $(\sqrt6-\sqrt2,\sqrt6+\sqrt2)$
\end{itemize}

\gkpart{解答题}
\begin{itemize}[leftmargin=2.6em]
  \item[\textbf{17.}] (1) $a_n=2n+1$；(2) $\frac{n}{3(2n+3)}$
  \item[\textbf{18.}] (1) 两平面垂直，见证明；(2) $\frac{\sqrt3}{3}$
  \item[\textbf{19.}] (1) $y=c+d\sqrt x$；(2) $\widehat y=100.6+68\sqrt x$；(3) ① $576.6$ t、$66.32$ 千元；② $x=46.24$ 千元（最大预报利润 $66.36$ 千元）
  \item[\textbf{20.}] (1) $y=\sqrt a\,x-a$ 与 $y=-\sqrt a\,x-a$；(2) 存在，$P=(0,-a)$
  \item[\textbf{21.}] (1) $a=-\frac34$；(2) 零点数为 $1$（$a<-\frac54$ 或 $a>-\frac34$），$2$（$a=-\frac54$ 或 $a=-\frac34$），$3$（$-\frac54<a<-\frac34$）
  \item[\textbf{22.}] (1) $DE$ 是圆的切线，见证明；(2) $60^\circ$
  \item[\textbf{23.}] (1) $C_1:\rho\cos\theta=-2$；$C_2:\rho^2-2\rho\cos\theta-4\rho\sin\theta+4=0$；(2) U：原文未定义三角形顶点 $C_2$；若指圆心则面积为 $\frac12$（条件性结论）
  \item[\textbf{24.}] (1) $(\frac23,2)$；(2) $a>2$
\end{itemize}
